% Figura: teorema de Stokes. La circulación en los bordes interiores de los
% parches infinitesimales se cancela y solo queda la del borde L = dS.
\begin{figure}[ht]
\centering
\begin{tikzpicture}[>=Stealth, x=1cm, y=1cm]
  \def\blob{plot[smooth cycle, tension=0.75] coordinates {(-4.3,0.1) (-3.1,1.25) (-1.2,1.5)
     (0.9,1.05) (2.7,1.45) (4.1,0.8) (4.5,-0.25) (3.3,-1.0) (1.0,-1.3) (-1.4,-1.1) (-3.3,-0.9)}}
  \path[name path=superficie] \blob;
  % superficie con degradado azul
  \shade[left color=azulclaro!35, right color=azulclaro!85] \blob;
  \begin{scope}
    \clip \blob;
    % parches con su circulación (elipses por la perspectiva)
    \foreach \j in {-3,-2,...,3}{\foreach \i in {-8,-7,...,8}{%
      \pgfmathsetmacro{\cx}{0.55*\i+0.12*\j}\pgfmathsetmacro{\cy}{0.33*\j-0.03*\i*\i/8}%
      \draw[black!70, line width=0.55pt,
            decoration={markings, mark=at position 0.2 with {\arrow{Stealth[length=2.6pt]}},
                                  mark=at position 0.7 with {\arrow{Stealth[length=2.6pt]}}},
            postaction={decorate}] (\cx,\cy) ellipse (0.27 and 0.145);}}
  \end{scope}
  % borde L con orientación
  \draw[rojoteo, line width=1.5pt,
        decoration={markings, mark=between positions 0.04 and 0.97 step 0.1 with {\arrow{Stealth[length=6pt]}}},
        postaction={decorate}] \blob;
  \draw[azulacero, line width=0.5pt, densely dashed] (0,0.1) -- ++(0,0.05);
  \draw[->, line width=1.6pt, azulacero] (0.15,0.25) -- ++(0,1.9) node[above, text=azulacero] {$\dS$};
  \node[font=\Large, text=azulprof] at (-3.4,0.1) {$S$};
  \node[font=\Large, text=rojoteo, anchor=west] at (4.6,-0.55) {$L=\partial S$};
\end{tikzpicture}
\caption{Idea del teorema de Stokes. La superficie $S$ se divide en parches
infinitesimales orientados, cada uno con su circulación (elipses). En las
paredes interiores cada lado se recorre dos veces en sentidos opuestos y la
circulación se cancela; solo sobrevive la circulación sobre el borde exterior
$L=\partial S$ (en rojo). La normal $\dS$ y el sentido de recorrido de $L$ están
ligados por la regla de la mano derecha.}
\label{fig:stokes}
\end{figure}
